% META: {"id": "TEXP-MAT-CR-010", "chapitre": "TEXP-MATRICES-MARKOV", "type_objet": "cours", "capacites_codes": ["C1"], "sources_inspiration": [], "mode_creation": "ex_nihilo", "status": "generated"}

\section{Opérations sur les matrices}

% ============================================================
% STRATE 1 — Definitions et operations de base
% ============================================================

\definition[1]{
Une \textbf{matrice} de taille $p\times q$ est un tableau de nombres réels comportant $p$ lignes et $q$ colonnes. Une matrice \textbf{carrée} a autant de lignes que de colonnes ($p=q$) ; une matrice \textbf{colonne} a une seule colonne ($q=1$) ; une matrice \textbf{ligne} a une seule ligne ($p=1$).
}

\propriete[Somme et produit]{
Deux matrices de même taille $p\times q$ s'\textbf{additionnent} coefficient par coefficient. Le \textbf{produit} $AB$ de deux matrices n'est défini que si le nombre de colonnes de $A$ égale le nombre de lignes de $B$ ; le coefficient en ligne $i$, colonne $j$ de $AB$ est la somme des produits des coefficients de la ligne $i$ de $A$ par ceux de la colonne $j$ de $B$.
}

\exemple{
Avec $A=\begin{pmatrix}2&1\\1&1\end{pmatrix}$ et $B=\begin{pmatrix}1&0\\3&2\end{pmatrix}$ :
$$A+B=\begin{pmatrix}3&1\\4&3\end{pmatrix}\qquad AB=\begin{pmatrix}5&2\\4&2\end{pmatrix}$$
}

% BEGIN-VERIFY
% from sympy import Matrix
% A = Matrix([[2,1],[1,1]])
% B = Matrix([[1,0],[3,2]])
% assert A + B == Matrix([[3,1],[4,3]])
% assert A * B == Matrix([[5,2],[4,2]])
% END-VERIFY

\erreurFrequente{
\textbf{Le produit matriciel n'est pas commutatif.} En général, $AB
eq BA$ (et l'un des deux produits peut même ne pas être défini si les tailles ne correspondent pas). Ne jamais échanger l'ordre des facteurs sans justification.
}

% ============================================================
% STRATE 2 — Inverse et puissances
% ============================================================

\definition[2]{
Une matrice carrée $A$ est \textbf{inversible} s'il existe une matrice $A^{-1}$, appelée \textbf{inverse} de $A$, telle que $AA^{-1}=A^{-1}A=I$ (matrice identité). Pour une matrice $2\times2$, $A=\begin{pmatrix}a&b\\c&d\end{pmatrix}$ est inversible si et seulement si son \textbf{déterminant} $\det(A)=ad-bc$ est non nul, et alors
$$A^{-1}=\dfrac{1}{ad-bc}\begin{pmatrix}d&-b\\-c&a\end{pmatrix}$$
}

\exemple{
Pour $A=\begin{pmatrix}2&1\\1&1\end{pmatrix}$, $\det(A)=2\times1-1\times1=1
eq0$ : $A$ est inversible, et
$$A^{-1}=\dfrac{1}{1}\begin{pmatrix}1&-1\\-1&2\end{pmatrix}=\begin{pmatrix}1&-1\\-1&2\end{pmatrix}$$
On vérifie que $AA^{-1}=I$.
}

% BEGIN-VERIFY
% from sympy import Matrix, eye
% A = Matrix([[2,1],[1,1]])
% assert A.det() == 1
% Ainv = A.inv()
% assert Ainv == Matrix([[1,-1],[-1,2]])
% assert A * Ainv == eye(2)
% END-VERIFY

\propriete[Puissances d'une matrice carrée]{
Pour une matrice carrée $A$ et un entier $n\geqslant1$, $A^n$ désigne le produit de $A$ par elle-même $n$ fois ($A^0=I$ par convention). Le calcul de $A^n$ se fait par produits matriciels successifs, ou (pour des valeurs de $n$ élevées) à l'aide d'un outil de calcul.
}

\exemple{
Pour $A=\begin{pmatrix}2&1\\1&1\end{pmatrix}$ : $A^2=\begin{pmatrix}5&3\\3&2\end{pmatrix}$ et $A^3=\begin{pmatrix}13&8\\8&5\end{pmatrix}$.
}

% BEGIN-VERIFY
% from sympy import Matrix
% A = Matrix([[2,1],[1,1]])
% assert A**2 == Matrix([[5,3],[3,2]])
% assert A**3 == Matrix([[13,8],[8,5]])
% END-VERIFY
